\documentclass[10pt,twocolumn]{article}
\usepackage[margin=1in]{geometry}
\usepackage{booktabs}
\usepackage{hyperref}
\usepackage{graphicx}
\usepackage{amsmath}
\usepackage{parskip}
\usepackage{microtype}
\usepackage{xcolor}
\usepackage{enumitem}
\usepackage{multirow}
\usepackage{tabularx}
\setlist{noitemsep}

\hypersetup{colorlinks=true, linkcolor=blue, citecolor=blue, urlcolor=blue}

\title{\textbf{Cross-City Urban Flood Detection from SAR:\\
       Geographic Context Features and the Data-Volume Ceiling}}
\author{Romit Basak\\
  \small CS5330 Pattern Recognition and Computer Vision --- Spring 2026\\
  \small Khoury College of Computer Sciences, Northeastern University\\
  \small Advisor: Prof.\ Bruce Maxwell}
\date{}

\begin{document}
\maketitle

% ============================================================
\begin{abstract}
Urban flood detection from Synthetic Aperture Radar (SAR) imagery is
critical for cloud-penetrating disaster response, but generalizing to
unseen cities remains an open problem.  We investigate two strategies
for improving cross-city generalization on UrbanSARFloods under a strict
event-held-out evaluation split: (1) static geographic context
features---urban morphology (WorldCover, Google Open Buildings),
topographic position (HAND elevation), and sensor geometry
(Sentinel-1 incidence angle); and (2) pseudo-label pretraining on
11 new flood events using AWEI\_sh optical confidence maps.
Our best model, an ImageNet-pretrained ResNet-34 U-Net fine-tuned
with focal loss ($\alpha_{\text{flood}}=0.75$, $\gamma=2$) on 7
training cities, achieves peak flood F1 = 0.150 and macro F1 = 0.551
on 7 completely unseen test cities.
All experimental conditions---SAR-only, geographic context augmentation,
and pseudo-label pretraining---converge to the same flood F1 ceiling of
$0.10$--$0.15$ (macro F1 $0.49$--$0.55$), with within-run variance
$\sigma \approx 0.020$ making differences between configurations
inconclusive.
We document four implementation pitfalls: a degenerate multi-class label
scheme, collapse-to-non-flood under class imbalance, Sen1Floods11
pretraining collapse, and pseudo-label pretraining collapse via rapid
overfitting to noisy optical-derived labels.
These findings collectively characterize what is tractable in cross-city
SAR urban flood detection with limited labeled data.
\end{abstract}

% ============================================================
\section{Introduction}

Flooding is the most economically damaging natural disaster globally,
affecting an estimated 1.8 billion people annually.  SAR satellites
such as Sentinel-1 provide cloud-penetrating flood maps within hours
of an overpass, making them the primary data source for operational
disaster response through services such as Copernicus EMS.

SAR flood detection in open areas is tractable: calm water produces
specular reflection that appears dark in SAR imagery against rougher
terrain.  Urban environments fundamentally break this assumption.
Building fa\c{c}ades produce double-bounce reflections that brighten
flooded streets.  Radar shadow behind tall buildings creates dark
regions that mimic open water.  Smooth impervious surfaces---roads,
parking lots---also appear dark even when dry.  As a result, a model
trained on one city's SAR flood signatures often fails on another city
with different building morphology, street geometry, and acquisition
geometry.

We investigate two hypotheses.  First, a model that knows
\emph{where it is}---what the urban density, topographic position
relative to drainage, and sensor viewing angle are at each
pixel---can reason from physics-grounded priors that transfer across
cities.  Second, pseudo-labeled chips from new flood events, derived
from AWEI\_sh optical water indices with confidence weighting, can
expand the effective training set and raise the data-volume ceiling.

\paragraph{Contributions.}
\begin{enumerate}
  \item A reproducible binary-label pipeline for cross-city SAR flood
        detection with strict event-held-out evaluation.
  \item Empirical comparison of SAR-only vs.\ geographic context
        augmented models (WorldCover + Google Buildings + HAND +
        incidence angle).
  \item An end-to-end AWEI\_sh pseudo-label pipeline evaluated on
        11 new flood events, with confidence-weighted focal loss
        and empirical evaluation of transfer to unseen cities.
  \item Diagnosis and documentation of four critical implementation
        failure modes with verified fixes.
  \item Characterization of a data-volume ceiling at flood F1
        $\approx 0.10$--$0.15$ (macro F1 $\approx 0.49$--$0.55$)
        that persists across all experimental conditions.
\end{enumerate}

% ============================================================
\section{Related Work}

\paragraph{SAR flood datasets.}
Sen1Floods11~\cite{bonafilia2020} provides 446 hand-labeled Sentinel-1
chips across 11 open-area flood events on 6 continents.
UrbanSARFloods~\cite{zhao2024} extends this to 8,879 chips from 18
urban flood events, but published baselines use random chip splits that
permit spatial autocorrelation leakage between training and test chips
from the same city.
Kuro Siwo~\cite{yadav2024} scales to 33 billion m$^2$ of SAR imagery
but focuses on open-area detection, explicitly noting cross-city
generalization as an unsolved problem.

\paragraph{Evaluation methodology.}
Al Mehedi et al.~\cite{almehedi2025} compare Otsu thresholding,
change detection, and Random Forest on Houston Harvey, reporting urban
flood F1 = 0.608 (ML classifier) with precision and recall
$\approx 0.60$.  However, their model is trained and tested on the
same event.  Zhao et al.\ report F1 = 0.51--0.77 with random chip
splits from the same events.  Both evaluations are in-domain.  Our
event-held-out split is strictly harder: no chip from any test city
appears during training.

\paragraph{Geographic priors for flood detection.}
HAND (Height Above Nearest Drainage)~\cite{renno2008} is a physically
motivated prior: pixels near drainage networks flood before those at
higher elevations, reflecting hydraulic principles that transfer across
cities rather than city-specific appearance.
DeepSARFlood (2025) uses HAND and terrain slope as both inputs and
postprocessors.  Sentinel-1 incidence angle effects on urban SAR
scattering are well documented~\cite{moreira2013} but rarely used as
explicit model inputs.

\paragraph{Optical pseudo-labels for SAR.}
AWEI\_sh (Automated Water Extraction Index with shadow
suppression)~\cite{feyisa2014} reduces building-shadow
misclassification in urban optical imagery relative to NDWI.
Using optical water indices to generate SAR training labels has been
proposed for open-area floods~\cite{bonafilia2020} but its
effectiveness in dense urban scenes at 10\,m resolution is not
established prior to this work.

% ============================================================
\section{Data}

\subsection{Datasets}

\textbf{UrbanSARFloods}~\cite{zhao2024} provides 8,879 Sentinel-1
chips from 18 urban flood events across 5 continents, organized into
\texttt{01\_NF} (confirmed non-flooded), \texttt{02\_FO} (flooded
open area), and \texttt{03\_FU} (flooded urban).

\textbf{Pseudo-labeled events.}  We construct a supplementary
pseudo-labeled dataset from 11 new flood events not present in
UrbanSARFloods: Brisbane 2022, Sydney 2022, Valencia 2024,
Liège 2021 (two SAR acquisition dates), Cologne 2021,
Zhengzhou 2021, Derna 2023 (two SAR dates), Emilia-Romagna 2023,
and Rio Grande do Sul 2024.
Events were identified via a Sentinel-1/Sentinel-2 co-registration
discovery pipeline requiring cloud cover $< 30$\% within $\pm$2 days
of the SAR acquisition.
Per-pixel confidence maps are derived from AWEI\_sh on
Sentinel-2 Surface Reflectance imagery, with shadow masking
(NIR $< 0.10$ and Blue $< 0.05$) and SCL-based cloud exclusion.
After spatial chipping (512$\times$512 at 10\,m, 50\% overlap stride)
and a 3:1 non-flood:flood subsampling cap, this yields 2,073 chips
from 11 events.

\subsection{Label Scheme}

We use a binary label scheme throughout: $1 =$ flood, $0 =$ non-flood,
$-1 =$ ignore.

\begin{itemize}
  \item \texttt{01\_NF} chips: all pixels $\to 0$ (confirmed non-flood).
  \item \texttt{02\_FO / 03\_FU}: background $\to -1$; open and urban
        flood $\to 1$.
  \item Pseudo-labeled chips: AWEI\_sh confidence $\geq 0.65 \to 1$
        (flood, loss weight $=$ confidence);
        confidence $\leq 0.35 \to 0$
        (non-flood, weight $= 1 -$ confidence);
        otherwise $\to -1$ (ignore).
\end{itemize}

\paragraph{Why binary?}
Any multi-class scheme where non-flood is one of the averaged classes
produces misleading macro metrics: non-flood pixels dominate ($>98$\%),
so a model predicting non-flood everywhere achieves macro F1 $> 0.35$
while flood F1 $\approx 0$.  We verified this empirically
(Section~\ref{sec:degenerate}).  We report \textbf{flood F1}
(class 1 only) as our primary operational metric and
\textbf{macro F1} (average of flood and non-flood F1) as a
secondary metric for comparability with standard CV benchmarks.
The all-non-flood degenerate baseline achieves macro F1 = 0.33.

\subsection{Event-Level Split}

\begin{table}[h]
\centering\small
\caption{Event-level train/test split (3,261 train / 3,210 test chips).
  No chip from any test city appears during training.
  Pseudo-labeled events are disjoint from both splits.}
\label{tab:split}
\begin{tabularx}{\columnwidth}{lX}
\toprule
Split & Events \\
\midrule
Train (7) & Houston, Beira, Japan, Canada, Iran, Lumberton, Somalia \\
Test  (7) & Hagibis, Sydney, Coraki, Niger, Hebei, Beledweyne, PortMacquarie \\
\bottomrule
\end{tabularx}
\end{table}

\subsection{Input Features}

\paragraph{SAR channels (4).}
VV intensity, VH intensity, local variance of VV ($5{\times}5$ window),
local variance of VH.  Variance channels disambiguate radar shadow
(near-zero variance) from open water (non-zero speckle variance).

\paragraph{Geographic context channels (4, aux models).}
(1) Google Open Buildings density (normalised to $[0,1]$ at 100\,m);
(2) ESA WorldCover v200 built-up fraction at 100\,m;
(3) HAND elevation clipped to 30\,m and normalised to $[0,1]$;
(4) Sentinel-1 incidence angle ($[20^\circ, 46^\circ] \to [0,1]$).
All exported from Google Earth Engine and resampled to chip resolution
via bilinear interpolation.  The same pipeline was applied to
pseudo-labeled event AOIs to produce 8-channel pseudo chips.

% ============================================================
\section{Method}

\subsection{Model Architecture}

We use a U-Net with ResNet-34 encoder from
\texttt{segmentation-models-pytorch}~\cite{iakubovskii2019},
pretrained on ImageNet.  The first convolutional layer is extended to
accept 8 input channels: weights for the original 3 channels are
preserved; additional channels are initialised as their mean, avoiding
large initial gradients.  The decoder predicts 2 classes (non-flood,
flood).

\subsection{Loss Function}

Standard cross-entropy and Dice losses collapse to the
non-flood-everywhere solution when flood pixels constitute $<1$\% of
labeled pixels.  We use focal loss~\cite{lin2017focal}:

\[
  \mathcal{L} = -\alpha_t (1 - p_t)^\gamma \log p_t, \quad \gamma = 2.0
\]

with $\alpha_{\text{flood}} = \text{fw}/(1+\text{fw})$.
For pseudo-labeled chips, per-pixel focal loss is additionally
multiplied by the AWEI\_sh confidence weight, so uncertain pixels
contribute proportionally less gradient.

\subsection{Training Protocol}

\paragraph{Standard fine-tuning.}
Up to 40 epochs, discriminative learning rates (encoder
$2{\times}10^{-5}$, decoder/head $10^{-4}$), cosine decay with 3-epoch
linear warmup, early stopping patience 8 on flood F1, batch size 16,
weight decay $5{\times}10^{-4}$, AdamW.

\paragraph{Pseudo-label pipeline.}
Stage~0 pretrains on 2,073 pseudo-labeled chips for up to 20 epochs
($\text{lr} = 10^{-4}$, patience 6 on zero-shot flood F1 monitored
on the UrbanSARFloods test set without fine-tuning).
Stage~1 fine-tunes on UrbanSARFloods using the same protocol above.

% ============================================================
\section{Results}

\subsection{Main Results}
\label{sec:results}

Table~\ref{tab:main} summarises all model configurations.
We report peak flood F1 (best single epoch), macro F1 at that epoch,
and mean$\,\pm\,$SD of flood F1 over all epochs.
Within-run SD $\approx 0.020$ is comparable to differences between
configurations, making most pairwise differences inconclusive.

% Full-width table to avoid column overflow
\begin{table*}[t]
\centering\small
\caption{Cross-city results on event-held-out test split (3,210 chips, 7 unseen cities).
  Macro F1 reported at the peak-flood-F1 epoch.
  Degenerate (all-non-flood) baseline: macro F1 = 0.33, flood F1 = 0.}
\label{tab:main}
\begin{tabular}{lcccc}
\toprule
Model & Channels & Peak Flood F1 & Macro F1 & Mean$\pm$SD Flood F1 \\
\midrule
SAR-only (fw=10)               & 4 & 0.120 & $\sim$0.48 & $0.082\pm0.020$ \\
Aux features (fw=10)           & 8 & 0.102 & $\sim$0.49 & $0.086\pm0.014$ \\
Aux features (fw=3)            & 8 & 0.150 & $\sim$0.52 & $0.100\pm0.022$ \\
\midrule
Finetune-only, fw=3$^\dagger$  & 8 & 0.143 & 0.551      & $0.103\pm0.019$ \\
Pseudo+Finetune, fw=3$^\dagger$& 8 & 0.136 & 0.550      & $0.091\pm0.025$ \\
\bottomrule
\multicolumn{5}{l}{$^\dagger$Macro F1 logged exactly from same codebase.
  Top three rows: macro F1 estimated from macro IoU.}
\end{tabular}
\end{table*}

All five configurations converge to flood F1 $0.10$--$0.15$ and
macro F1 $0.49$--$0.55$.
Finetune-only (0.143) and Aux fw=3 (0.150) differ by less than one
within-run SD, confirming equivalence.  Pseudo-label pretraining
(0.136) also falls within this band.  No configuration meaningfully
outperforms any other.

\subsection{Per-City Breakdown}

Table~\ref{tab:percity} shows per-city flood F1, precision, and recall
for the SAR-only and best Aux configurations at their best checkpoints.

% Single-column with abbreviated names and tighter spacing
\begin{table}[h]
\centering\small
\setlength{\tabcolsep}{4pt}
\caption{Per-city flood F1 at best checkpoint (P = precision, R = recall).
  Aux = 8-channel geographic context, fw=3.}
\label{tab:percity}
\begin{tabular}{lc rrr rrr}
\toprule
 & & \multicolumn{3}{c}{SAR-only (fw=10)} & \multicolumn{3}{c}{Aux (fw=3)} \\
\cmidrule(lr){3-5}\cmidrule(lr){6-8}
City & $N$ & F1 & P & R & F1 & P & R \\
\midrule
Port Mac. & 187 & \textbf{.541} & .435 & .716 & .200 & .549 & .122 \\
Coraki    & 135 & .282 & .167 & .905 & \textbf{.420} & .431 & .410 \\
Niger     & 501 & .147 & .081 & .792 & .151 & .108 & .251 \\
Beledw.   & 492 & .092 & .050 & .532 & \textbf{.170} & .094 & .899 \\
Sydney    & 962 & .091 & .051 & .424 & \textbf{.106} & .064 & .311 \\
Hebei     & 933 & .054 & .028 & .512 & \textbf{.125} & .071 & .518 \\
\midrule
Mean      &3210 & .201 & & & .195 & & \\
\bottomrule
\end{tabular}
\end{table}

SAR-only (fw=10) achieves high recall (0.42--0.91) at very low
precision (0.03--0.44), reflecting systematic over-prediction of flood.
Aux fw=3 rebalances to higher precision (0.06--0.55) and lower recall
(0.12--0.90), producing higher F1 for 5 of 6 cities.

PortMacquarie and Coraki (Australian river floods, large contiguous
water bodies) are easiest.  Hebei (dense urban China) and Sydney
(fragmented urban flooding) are hardest: fw=3 still reaches only
F1 = 0.125 and 0.106 respectively.

\subsection{Comparison with Published Results}

\begin{table}[h]
\centering\small
\caption{Comparison with published methods.  Direct numerical
  comparison is not valid due to evaluation protocol differences.}
\label{tab:compare}
\begin{tabular}{lcc}
\toprule
Method & F1 & Evaluation \\
\midrule
Al Mehedi RF~\cite{almehedi2025}   & 0.608      & Same event \\
Zhao et al.~\cite{zhao2024}        & 0.51--0.77 & Random chip split \\
Ours                               & 0.150      & Cross-city held-out \\
\bottomrule
\end{tabular}
\end{table}

The gap between published F1 $\approx 0.6$ and our flood F1 $\approx
0.15$ quantifies the cost of cross-city generalization.  On macro F1
our models (0.49--0.55) substantially exceed the all-non-flood baseline
(0.33), confirming genuine learning rather than degenerate prediction.

% ============================================================
\section{Failure Mode Analysis}
\label{sec:failure}

\subsection{Degenerate Multi-Class Labeling}
\label{sec:degenerate}

Our initial 4-class scheme (0=ignore, 1=open flood, 2=urban flood,
3=non-flood) reported macro F1 up to 0.389.  Diagnosis revealed the
model predicted non-flood for 88.5\% of pixels; true flood F1 was
0.116.  Non-flood dominates the dataset, so near-perfect non-flood F1
inflates the macro average regardless of flood detection quality.
The scheme also introduced a cross-dataset label conflict between
Sen1Floods11 and UrbanSARFloods class indices.

\textbf{Fix:} Binary labels; flood F1 as primary metric.

\subsection{Collapse-to-Non-Flood}
\label{sec:collapse}

Cross-entropy and Dice losses both collapsed to all-non-flood by epoch
3.  Dice loss finds a local minimum at zero flood predictions because
its denominator is dominated by non-flood volume at $<1$\% flood pixel
prevalence.  Focal loss with $\alpha_{\text{flood}}=0.75$, $\gamma=2.0$
resolved the collapse by down-weighting easy non-flood examples.

\subsection{Sen1Floods11 Pretraining Collapse}
\label{sec:s1collapse}

Pretraining on Sen1Floods11 collapsed after epoch 2: flood F1 fell from
0.289 to 0.063 despite decreasing training loss.  The dataset has
many all-NaN chips and no dense urban scenes.  We use ImageNet
initialization directly for all reported models.

\subsection{Pseudo-Label Pretraining Collapse}
\label{sec:pseudo_collapse}

Pseudo-label pretraining on 2,073 chips from 11 events degraded peak
flood F1 from 0.143 to 0.136 (finetune-only baseline).
Table~\ref{tab:pseudo_curve} shows the Stage~0 curve: zero-shot F1
peaked at 0.124 at epoch~3, then collapsed to near zero by epoch~9.

\begin{table}[h]
\centering\small
\caption{Stage~0 pseudo-pretraining log.  Zero-shot flood F1 on the
  UrbanSARFloods test set (no fine-tuning).}
\label{tab:pseudo_curve}
\begin{tabular}{rcc}
\toprule
Epoch & Train loss & Zero-shot F1 \\
\midrule
1 & 0.0462 & 0.033 \\
2 & 0.0101 & 0.025 \\
3 & 0.0065 & \textbf{0.124} \\
4 & 0.0043 & 0.051 \\
5 & 0.0040 & 0.052 \\
6 & 0.0034 & 0.013 \\
7 & 0.0031 & 0.010 \\
8 & 0.0030 & 0.013 \\
9 & 0.0029 & 0.001\quad$\leftarrow$stopped \\
\bottomrule
\end{tabular}
\end{table}

Three factors explain the collapse.  \textbf{Dataset size}: 2,073
chips is insufficient to stabilize the encoder.  \textbf{Label
sparsity}: four of eleven events (Cologne, Zhengzhou, Liège~B,
Emilia-Romagna) produced zero flood chips because AWEI\_sh cannot
detect flood water in dense urban cores where building shadows suppress
the optical water signal; only 621 of 2,073 chips contained any flood
signal.  \textbf{Domain mismatch}: pseudo-labeled events are
predominantly open fluvial floods, not adding the dense urban SAR
diversity needed to improve on Hebei and Sydney.

\textbf{Implication:} SAR change detection (pre-flood vs.\ flood-period
backscatter ratio) is a more promising pseudo-label strategy, as it
avoids optical cloud and shadow entirely.

% ============================================================
\section{Discussion}

\subsection{The Data-Volume Ceiling}

All experimental conditions converge to flood F1 $\approx 0.10$--$0.15$
and macro F1 $\approx 0.49$--$0.55$.  This ceiling persists across
feature sets, loss weight tuning, and pseudo-label pretraining.  It
is not a model capacity problem (training loss $\to 0.007$) nor a label
quality problem.  The pseudo-label experiment demonstrates it is not
solely a data quantity problem---label quality and urban diversity of
pseudo-labeled events are equally binding constraints.

Kuro Siwo~\cite{yadav2024} identified cross-city generalization as
unsolved even at 33 billion m$^2$.  Our results provide a concrete
ceiling characterization at 7 labeled training cities with two
negative results: geographic context features and optical
pseudo-labeling do not overcome it at this scale.

\subsection{Why Geographic Features Did Not Help}

The per-city results show geographic context helps where it should:
Coraki and Beledweyne improve (0.282$\to$0.420 and 0.092$\to$0.170),
both events where flood extent correlates with topographic position.
Hebei and Sydney show negligible improvement, consistent with dense
urban scattering complexity exceeding what 7 training cities can teach.
At 20--30 training cities, geographic context should show clear benefit.

\subsection{Implications for Kolkata Deployment}

Our results suggest: (1) hand-labeled South and Southeast Asian urban
flood events are essential; (2) SAR change detection pseudo-labels
are more promising than AWEI\_sh; (3) geographic context features
are worth retaining for when sufficient training cities are available.

% ============================================================
\section{Conclusion}

We presented a systematic investigation of cross-city SAR urban flood
detection across five experimental conditions.  Our best model achieves
peak flood F1 = 0.150, macro F1 = 0.551 on 7 unseen cities.
Geographic context features and AWEI\_sh pseudo-label pretraining
on 11 additional events neither improved nor degraded results beyond
the within-run noise floor, confirming a data-volume ceiling that
persists across all conditions tested.

Four documented failure modes provide concrete guidance for future work.
The path to substantially higher cross-city flood F1 runs through
more hand-labeled urban events from diverse morphological contexts.

Future directions: SAR change detection for pseudo-label generation;
geospatial foundation model encoders (Prithvi-EO-2.0); HAND-based
flood shadow interpolation; scaling to 20+ labeled training cities.

% ============================================================
\bibliographystyle{plain}
\begin{thebibliography}{9}

\bibitem{bonafilia2020}
D.~Bonafilia, B.~Tellman, T.~Anderson, and E.~Issenberg,
``Sen1Floods11: A georeferenced dataset to train and test deep learning
flood algorithms for Sentinel-1,'' in \emph{CVPR Workshops}, 2020.

\bibitem{zhao2024}
B.~Zhao, H.~Sui, and J.~Liu,
``UrbanSARFloods: Sentinel-1 SLC-based benchmark dataset for urban and
open-area flood mapping,'' in \emph{CVPR Workshop on Earth Vision}, 2024.

\bibitem{almehedi2025}
M.~A.~Al Mehedi, V.~Smith, and P.~Kremer,
``A comparative analysis of urban and peri-urban flood identification
using SAR imagery,'' \emph{PLOS Water}, vol.~4, no.~9, p.~e0000269,
2025.

\bibitem{renno2008}
C.~D.~Renn\'{o} et al.,
``HAND, a new terrain descriptor using SRTM-DEM: Validation in the
Amazon,'' \emph{Remote Sensing of Environment}, vol.~112, no.~9, 2008.

\bibitem{lin2017focal}
T.-Y.~Lin, P.~Goyal, R.~Girshick, K.~He, and P.~Doll\'{a}r,
``Focal loss for dense object detection,'' in \emph{ICCV}, 2017.

\bibitem{yadav2024}
R.~Yadav et al.,
``Kuro Siwo: 33 billion m$^2$ under the water --- A global multi-temporal
SAR dataset,'' in \emph{NeurIPS Datasets and Benchmarks Track}, 2024.

\bibitem{iakubovskii2019}
P.~Iakubovskii,
``Segmentation Models PyTorch,'' GitHub, 2019.
\url{https://github.com/qubvel/segmentation_models.pytorch}

\bibitem{moreira2013}
A.~Moreira et al.,
``A tutorial on synthetic aperture radar,''
\emph{IEEE Geoscience and Remote Sensing Magazine}, vol.~1, no.~1, 2013.

\bibitem{feyisa2014}
G.~L.~Feyisa, H.~Meilby, R.~Fensholt, and S.~R.~Proud,
``Automated Water Extraction Index: A new technique for surface water
mapping using Landsat imagery,'' \emph{Remote Sensing of Environment},
vol.~140, pp.~23--35, 2014.

\end{thebibliography}

\end{document}
